\documentclass[aspectratio=169]{beamer}

\title{GCC Record: Embedding Build Inputs in ELF Binaries}
\subtitle{A compiler and linker plugin experiment in build provenance}
\date{FOSDEM 2027 --- GCC/GNU Toolchain devroom}

\begin{document}

\begin{frame}
  \titlepage
\end{frame}

\begin{frame}{The build-input visibility gap}
  \begin{itemize}
    \item A build consumes more than the source file named on the command line
    \item Compilation reads headers and other files; linking consumes objects, libraries, and startup files
    \item Once the build finishes, those inputs are not necessarily easy to discover from the output
    \item Separate logs and sidecar manifests can become detached from the artifact
  \end{itemize}
\end{frame}

\begin{frame}{Why this is useful}
  \begin{itemize}
    \item Investigate which inputs contributed to a build artifact
    \item Make build outputs easier to inspect during debugging and review
    \item Keep a compact input record alongside the object or executable
    \item Give downstream tooling a simple section to extract and analyze
    \item Hashes help distinguish recorded file contents; they are not a signature or proof that a build is trustworthy
  \end{itemize}
\end{frame}

\begin{frame}{Previous solutions: source identification keywords}
  \begin{itemize}
    \item SCCS and CVS supported identification keywords embedded in source text
    \item Examples include SCCS \texttt{\%W\%} / \texttt{\%G\%} and CVS \texttt{\$Id\$} / \texttt{\$Revision\$}
    \item These are useful for identifying individual source files and their revisions
    \item They depend on keywords being present in source and expanded by the version-control workflow
    \item They do not automatically provide a consolidated record of the files actually read across compilation and linking
  \end{itemize}
\end{frame}

\begin{frame}{What gcc-record does}
  \begin{itemize}
    \item Records the files observed during compilation and linking
    \item Stores a path and SHA-1 digest for each recorded file
    \item Embeds the record in a non-allocated, read-only ELF section named \texttt{files\_read}
    \item Supports plain text, gzip, and zstd section contents
    \item Provides a wrapper that loads the compiler and linker plugins automatically
  \end{itemize}
\end{frame}

\begin{frame}{A simple workflow}
  \begin{itemize}
    \item Build the plugins with \texttt{make}
    \item Compile and link using \texttt{gcc-record} in place of \texttt{gcc}:
      \begin{itemize}
	\item \texttt{./gcc-record -c hello.c -o hello.o}
	\item \texttt{./gcc-record -o hello hello.o}
      \end{itemize}
    \item Optional record compression: \texttt{--record=gz} or \texttt{--record=zstd}
    \item Other compiler arguments, including GCC's \texttt{-z} and \texttt{-Z}, are passed through
  \end{itemize}
\end{frame}

\begin{frame}{Compilation: record the translation unit}
  \begin{itemize}
    \item The GCC plugin observes included files and the main input file
    \item At the end of compilation, it creates the \texttt{files\_read} section in the object
    \item Each record is one path, a tab, and a SHA-1 digest per line
    \item Section flags make it read-only and not loaded as part of the program image
  \end{itemize}
\end{frame}

\begin{frame}{Linking: collect and merge}
  \begin{itemize}
    \item The linker plugin sees the linker's input files
    \item It records their paths and digests, and scans ELF inputs for existing \texttt{files\_read} sections
    \item Existing sections can be plain, gzip-compressed, or zstd-compressed
    \item The plugin combines records and removes duplicate lines
    \item The linker plugin writes a merged temporary record; the wrapper uses \texttt{objcopy} to install one final section
  \end{itemize}
\end{frame}

\begin{frame}{Inspecting the result}
  \begin{itemize}
    \item Confirm the section exists with \texttt{readelf -S -W hello}
    \item Extract it with:
      \begin{itemize}
	\item \small\texttt{objcopy --dump-section files\_read=inputs.txt hello /dev/null}
      \end{itemize}
    \item The extracted list can be reviewed or consumed by downstream scripts
    \item The record stays with the object or executable when copied as a file
  \end{itemize}
\end{frame}

\begin{frame}{Demonstration}
  \begin{itemize}
    \item Compile two source files separately with \texttt{gcc-record}
    \item Show the \texttt{files\_read} section in each object
    \item Link the objects and inspect the resulting executable's section
    \item Check that the final list includes source/header files, object files, and linker inputs
    \item Repeat with gzip and zstd modes; run the generated program
  \end{itemize}
\end{frame}

\begin{frame}{Scope and limitations}
  \begin{itemize}
    \item The current implementation targets ELF and GCC's plugin interface
    \item It records files observed by the compiler and linker plugins; it is not a full hermetic-build system
    \item Paths depend on the build environment and may not be portable
    \item Consumers should treat the section as build metadata, not as a security attestation
    \item Compression reduces section size but requires a compatible decoder to inspect the contents
  \end{itemize}
\end{frame}

\begin{frame}{Wrap-up and discussion}
  \begin{itemize}
    \item gcc-record puts a useful record of observed build inputs inside ELF artifacts
    \item Compiler and linker plugins capture per-object records and combine them at link time
    \item The result is directly inspectable with familiar binutils tools
    \item Discussion: what additional metadata or integration would make this more useful in real toolchains?
  \end{itemize}
\end{frame}

\end{document}
